\section*{Capítulo \ref{ch:b1:water-small-molecules} --- El agua y las biomoléculas pequeñas}

\begin{solution}{exo:b1:water-small-molecules:1}
La molécula, angular y \omterm{def:b1:water-small-molecules:water}{polar}, tiene un oxígeno parcialmente negativo e
hidrógenos parcialmente positivos: rodea un catión con sus oxígenos y un
anión con sus hidrógenos, y los estabiliza (hidratación); además forma
\omterm{def:b1:water-small-molecules:water}{enlaces de hidrógeno} con los grupos \omterm{def:b1:water-small-molecules:water}{polares}. El aceite no tiene cargas ni
grupos \omterm{def:b1:water-small-molecules:water}{polares} con los que enlazarse; las moléculas de agua que lo rodean
deben ordenarse en una jaula, lo cual es desfavorable, así que el aceite
queda excluido.
\end{solution}

\begin{solution}{exo:b1:water-small-molecules:2}
\omterm{def:b1:water-small-molecules:water}{Enlace de hidrógeno} (\qtyrange{4}{20}{kJ/mol}): apareamiento de \omterm{def:b1:water-small-molecules:ph}{bases} en el
ADN, estructura secundaria de las proteínas. \omterm{def:b1:water-small-molecules:bonds}{Enlace iónico} (unos
\qty{12}{kJ/mol} en agua): puentes salinos en las proteínas, fijación del
sustrato. \omterm{prop:b1:water-small-molecules:properties}{Efecto hidrófobo} (no es un enlace, pero es comparable en total):
ensamblaje de membranas, plegamiento de proteínas. \omterm{def:b1:water-small-molecules:bonds}{Van der Waals}
(\qtyrange{1}{4}{kJ/mol} cada una): el empaquetamiento estrecho de dos
superficies complementarias cualesquiera, el ajuste entre enzima y
sustrato.
\end{solution}

\begin{solution}{exo:b1:water-small-molecules:3}
$\mathrm{pH} = -\log(4\times 10^{-8}) = 7.4$. A \omterm{def:b1:water-small-molecules:ph}{pH} 1,5,
$[\mathrm{H^+}] = 10^{-1.5} = \qty{0.032}{mol/L}$.
\end{solution}

\begin{solution}{exo:b1:water-small-molecules:4}
De 5,8 a 7,8 aproximadamente (p$K_a \pm 1$). A 7,4:
$\log(\text{base}/\text{ácido})
= 0.6$, cociente $4$.
\end{solution}

\begin{solution}{exo:b1:water-small-molecules:5}
A \omterm{def:b1:water-small-molecules:ph}{pH} 4,76: el \qty{50}{\%}. A 7,2: $\log(\mathrm{A^-}/\mathrm{HA}) =
2.44$, cociente $275$: ionizado al \qty{99.6}{\%}. A \omterm{def:b1:water-small-molecules:ph}{pH} 2: cociente
$10^{-2.76} =
1/575$: ionizado al \qty{0.2}{\%}. La forma ácida neutra
atraviesa la \omterm{def:b1:membranes-transport:membrane}{bicapa lipídica}; el ácido acético se absorbe en el estómago,
donde no está ionizado, y queda atrapado como acetato en el \omterm{def:b1:eukaryotic-cell:organelle}{citosol}.
\end{solution}

\begin{solution}{exo:b1:water-small-molecules:6}
\omterm{def:b1:water-small-molecules:ph}{pH} $= 6.8 + \log 1 = 6.8$. Tras el HCl: \omterm{def:b1:water-small-molecules:ph}{base} $40$, ácido $60$: \omterm{def:b1:water-small-molecules:ph}{pH} $= 6.8 +
\log(40/60) = 6.62$. En agua pura, $[\mathrm{H^+}] = \qty{0.01}{mol/L}$:
\omterm{def:b1:water-small-molecules:ph}{pH} $2$.
\end{solution}

\begin{solution}{exo:b1:water-small-molecules:7}
$500/2.4 = \qty{208}{g}$. Para vaporizar una molécula hay que romper sus
cuatro \omterm{def:b1:water-small-molecules:water}{enlaces de hidrógeno}; \qty{2.4}{kJ/g} son \qty{43}{kJ/mol},
aproximadamente la energía de dos enlaces por molécula (cada enlace lo
comparten dos).
\end{solution}

\begin{solution}{exo:b1:water-small-molecules:8}
$\Delta G^{\circ\prime} = -RT\ln 19 = -2.48\times 2.94 =
\qty{-7.3}{kJ/mol}$. En la \omterm{def:b1:cell-unit-of-life:cell}{célula}: $\Delta G = -7.3 + 2.48\ln(2/0.02) =
-7.3 + 11.4 = \qty{+4.1}{kJ/mol}$: la reacción transcurre en sentido
inverso, hacia la glucosa-1-fosfato (como ocurre en la síntesis de
glucógeno).
\end{solution}

\begin{solution}{exo:b1:water-small-molecules:9}
El agua es más densa a \qty{4}{\celsius}; el agua más fría y el hielo son
más ligeros y se quedan en la superficie, de modo que el hielo se forma
arriba y aísla el líquido de debajo, que se mantiene cerca de
\qty{4}{\celsius} todo el invierno. Si el hielo se hundiera, los lagos se
congelarían por completo de abajo arriba y solo se deshelarían en la
superficie en verano; la vida de agua dulce no podría pasar un invierno.
\end{solution}

\begin{solution}{exo:b1:water-small-molecules:10}
$RT = \qty{2.58}{kJ/mol}$. En reposo: $Q = 0.9\times 10^{-3}\times 8\times
10^{-3}/8\times 10^{-3} = 9\times 10^{-4}$, $\ln Q = -7.0$: $\Delta G =
-30.5 - 18.1 = \qty{-48.6}{kJ/mol}$. Fatigado: $Q = 3\times 10^{-3}\times
30\times 10^{-3}/5\times 10^{-3} = 0.018$, $\ln Q = -4.0$: $\Delta G =
-30.5 - 10.4 = \qty{-40.9}{kJ/mol}$. La fatiga ha recortado la energía por
\omterm{def:b1:water-small-molecules:atp}{ATP} un \qty{16}{\%}, bastante para frenar las \omterm{def:b1:membranes-transport:transporters}{bombas} y los motores que más
la necesitan.
\end{solution}

\begin{solution}{exo:b1:water-small-molecules:11}
$[\mathrm{H^+}] = 10^{-7.2} = \qty{6.3e-8}{mol/L}$; volumen
\qty{1e-12}{L}: $\qty{6.3e-20}{mol}\times\num{6e23} = 38$ protones. Frente a
$\num{3e9}$ grupos tamponantes: el ``\omterm{def:b1:water-small-molecules:ph}{pH}'' no es un recuento de protones,
sino el estado de protonación de los tampones, que intercambian protones con
el agua mil millones de veces más deprisa de lo que los pocos protones
libres podrían importar.
\end{solution}

\begin{solution}{exo:b1:water-small-molecules:12}
En el equilibrio el \omterm{def:b1:water-small-molecules:atp}{ATP}, el ADP y el fosfato quedarían en un cociente
cercano a $10^{-5}$; la \omterm{def:b1:cell-unit-of-life:cell}{célula} mantiene el \omterm{def:b1:water-small-molecules:atp}{ATP} mil veces por encima del ADP,
de modo que su \omterm{prop:b1:water-small-molecules:families}{hidrólisis} libera \qty{50}{kJ/mol} en vez de nada. Todo
proceso cuesta arriba --- síntesis, transporte, movimiento --- se acopla a
ese desplazamiento, y el único papel del catabolismo es mantenerlo. Una
\omterm{def:b1:cell-unit-of-life:cell}{célula} cuyo \omterm{def:b1:water-small-molecules:atp}{ATP} alcanza el equilibrio no tiene gradiente que gastar: las
\omterm{def:b1:membranes-transport:transporters}{bombas} se paran, los iones se fugan, la \omterm{def:b1:cell-unit-of-life:cell}{célula} se hincha y muere en minutos.
Estar vivo es el mantenimiento de ese desequilibrio, no una sustancia.
\end{solution}

\begin{solution}{pb:b1:water-small-molecules:1}
\textbf{1.} $6.1 + \log(24/1.2) = 6.1 + 1.301 = 7.40$.
\textbf{2.} $10^{-7.4} = \qty{4.0e-8}{mol/L} = \qty{40}{nmol/L}$.
\textbf{3.} $40\times 5 = \qty{200}{nmol}$, seiscientas mil veces menos que
el bicarbonato.
\textbf{4.} $\mathrm{CO_2} = 1.2\times 46/40 = \qty{1.38}{mmol/L}$;
\omterm{def:b1:water-small-molecules:ph}{pH} $= 6.1 + \log(25/1.38) = 6.1 + 1.258 = 7.36$.
\textbf{5.} En un matraz, el cociente $20:1$ deja poco socio ácido para
absorber \omterm{def:b1:water-small-molecules:ph}{base}, y la capacidad es baja a una unidad de la p$K_a$; en el
cuerpo el socio ácido es un gas cuya concentración fijan los pulmones y que
pueden regenerar sin límite, de modo que el cociente puede reajustarse por
ventilación al margen de la p$K_a$.
\textbf{6.} $24/1.2 = 20$; la \omterm{def:b1:water-small-molecules:ph}{base} es el $20/21 = \qty{95}{\%}$ del par.
\textbf{7.} $60/5 = \qty{12}{mmol/L}$.
\textbf{8.} Bicarbonato $24 - 12 = \qty{12}{mmol/L}$; $\mathrm{CO_2}$,
$1.2 + 12 = \qty{13.2}{mmol/L}$.
\textbf{9.} \omterm{def:b1:water-small-molecules:ph}{pH} $= 6.1 + \log(12/13.2) = 6.1 - 0.04 = 6.06$.
\textbf{10.} $[\mathrm{H^+}] = \qty{0.012}{mol/L}$: \omterm{def:b1:water-small-molecules:ph}{pH} $1.92$.
\textbf{11.} Un desplazamiento de $1.34$ unidades: salva a la sangre del \omterm{def:b1:water-small-molecules:ph}{pH}
1,9, pero 6,06 queda muy por debajo de lo que tolera cualquier enzima. Malo,
como sistema cerrado.
\textbf{12.} El $[\mathrm{H^+}]$ pasa de \qty{40}{nmol/L} a $10^{-6.06} = \qty{870}{nmol/L}$: una subida de \qty{0.83}{\micro mol/L} para \qty{12}{mmol/L} añadidos, es decir, un protón de cada \num{14000} queda libre; el resto están sobre el bicarbonato.
\textbf{13.} \omterm{def:b1:water-small-molecules:ph}{pH} $= 6.1 + \log(12/1.2) = 6.1 + 1.0 = 7.10$.
\textbf{14.} El $\mathrm{CO_2}$ generado: $\qty{12}{mmol/L}\times
\qty{5}{L} = \qty{60}{mmol}$.
\textbf{15.} \qty{60}{mmol} en un minuto por encima de \qty{10}{mmol}:
siete veces más.
\textbf{16.} $\mathrm{CO_2} = \qty{0.9}{mmol/L}$; \omterm{def:b1:water-small-molecules:ph}{pH} $= 6.1 +
\log(12/0.9) = 6.1 + 1.125 = 7.22$.
\textbf{17.} Bicarbonato abierto cerca de 7,4: el \omterm{def:b1:water-small-molecules:ph}{pH} cambia una unidad
cuando el bicarbonato cae diez veces, de \qty{24}{} a \qty{2.4}{mmol/L}, es
decir, \qty{21.6}{mmol/L} de ácido por unidad; pero en las primeras unidades
la pendiente local es $2.3\times[\mathrm{HCO_3^-}] =
\qty{55}{mmol/L}$ por
unidad. Frente a los \qty{25}{mmol/L} por unidad de las proteínas, el
bicarbonato capta alrededor del $55/80 = \qty{69}{\%}$ de los protones, y
las proteínas el \qty{31}{\%}.
\textbf{18.} El bicarbonato capta $0.69\times 12 = \qty{8.3}{mmol/L}$:
$24 - 8.3 = \qty{15.7}{mmol/L}$; \omterm{def:b1:water-small-molecules:ph}{pH} $= 6.1 + \log(15.7/1.2) = 6.1 +
1.117 = 7.22$ (la parte de las proteínas es coherente con
\qty{25}{mmol/L} por unidad $\times 0.18$ unidades $= \qty{4.5}{mmol/L}$,
cerca de sus \qty{3.7}{mmol/L}).
\textbf{19.} Retirar el ácido regenera el bicarbonato consumido (de
\qty{12}{mmol/L} de vuelta a \qty{24}{}) y, con el $\mathrm{CO_2}$
mantenido, el \omterm{def:b1:water-small-molecules:ph}{pH} vuelve a 7,40.
\textbf{20.} $60/5 = \qty{12}{mmol/L}$ de bicarbonato perdidos al día; al
cabo de tres días, $24 - 36 < 0$: el bicarbonato se agota y el \omterm{def:b1:water-small-molecules:ph}{pH} se
desploma muy por debajo de 7 --- en la práctica hay que administrar
bicarbonato al paciente o dializarlo; tras un solo día,
$6.1 + \log(12/1.2) =
7.10$.
\textbf{21.} $\mathrm{CO_2} = 1.2\times 60/40 = \qty{1.8}{mmol/L}$;
\omterm{def:b1:water-small-molecules:ph}{pH} $= 6.1 + \log(24/1.8) = 6.1 + 1.125 = 7.22$. Con el bicarbonato a
\qty{33}{}: $6.1 + \log(33/1.8) = 6.1 + 1.263 = 7.36$.
\textbf{22.} Los pulmones cambian el denominador ($\mathrm{CO_2}$ disuelto)
en unas cuantas respiraciones, porque es un gas que exhalan; los riñones
cambian el numerador (el bicarbonato) excretando ácido y regenerando
bicarbonato, un proceso de horas a días.
\textbf{23.} A \omterm{def:b1:water-small-molecules:ph}{pH} 5, $[\mathrm{H^+}] = \qty{10}{\micro mol/L}$: \qty{15}{\micro mol} libres en \qty{1.5}{L}, una cuatromilésima de los \qty{60}{mmol}; el ácido sale unido a tampones, como amonio y como dihidrogenofosfato.
\textbf{24.} Perder ácido sube el \omterm{def:b1:water-small-molecules:ph}{pH}: el bicarbonato sube
$100/5 = \qty{20}{mmol/L}$ hasta \qty{44}{}; \omterm{def:b1:water-small-molecules:ph}{pH} $= 6.1 + \log(44/1.2) =
6.1 + 1.564 = 7.66$ (alcalosis), antes de que compensen los pulmones y los
riñones.
\textbf{25.} Cerrado: de 7,40 a 6,06, un desplazamiento de $-1.34$.
Abierto, solo con bicarbonato: hasta 7,10, un desplazamiento de $-0.30$.
Abierto con las proteínas: hasta 7,22, un desplazamiento de $-0.18$.
\end{solution}
